\begin{table}[ht]
\caption{\textbf{TheoremQA Questions example.} Both the query and the positive document uses the Pigeonhole Principle to solve the problem. The negative document is retrieved by BM25, and it only requires simple arithmetics to solve.}
\centering
% \resizebox{\textwidth}{!}{
\begin{tabular}{p{13cm}}
\toprule
% \hspace{5.5cm} \sustainable \\
% \midrule
\textit{\textbf{Query}} \\
\midrule
Mary is planning to bake exactly 10 cookies, and each cookie may be one of three different shapes -- triangle, circle, and square. Mary wants the cookie shapes to be a diverse as possible. What is the smallest possible count for the most common shape across the ten cookies?\\
\midrule
\textit{\textbf{Chain-of-thought reasoning to find documents}} \\
\midrule
This problem is similar to the formulation of Pigeonhole Principle. We can find solutions to other examples that can also be solved using this principle.\\
\midrule
\textit{\textbf{Example positive document}} \\
\midrule
Arbitrarily place 19 points in a unit square and cover as many of these points as possible with a circle of diameter $\frac{\sqrt 2}{3}$. Question: At least how many points can be guaranteed to be covered?

We can divide the unit square into 9 smaller squares, each with side length 1/3. Since there are 19 points and 9 smaller squares, by the Pigeonhole Principle, at least one of these smaller squares must contain at least 3 points.

Now, consider a circle with diameter $\frac{\sqrt 2}{3}$. The diagonal of each smaller square has length $\frac{\sqrt 2}{3}$, so the circle can be placed inside the smaller square in such a way that it covers the entire square. Therefore, the circle can cover at least 3 points in that square.

Therefore, the answer is 3.\\
\midrule
\textit{\textbf{Example negative document}} \\
\midrule
Yanna baked twenty butter cookies and forty biscuits in the morning.  In the afternoon, she baked ten butter cookies and twenty biscuits. How many more biscuits did she bake than butter cookies?
She baked 20 + 10 = 30 butter cookies.
And, she baked 40 + 20 = 60 biscuits.
Therefore, she baked 60 - 30 = 30 more biscuits than butter cookies.
The answer is 30
\\
\bottomrule
\end{tabular}
% }
\label{tab:example_theoremqa}
\end{table}